% Paquete institucional — autorización pedagógica (A1) + solicitud de fondos producción
% Distinto del PAQUETE-REMISION-CEI. Build: make institucional
\documentclass[11pt,a4paper]{scrartcl}

\usepackage[spanish]{babel}
\usepackage{amssymb}
\usepackage{microtype}
\usepackage{geometry}
\usepackage{booktabs}
\usepackage{setspace}
\IfFileExists{enumitem.sty}{%
  \usepackage{enumitem}%
  \setlist[itemize]{leftmargin=1.4em, itemsep=0.12em, topsep=0.15em}%
  \setlist[enumerate]{leftmargin=1.6em, itemsep=0.12em, topsep=0.15em}%
}{%
  \setlength{\leftmargini}{1.4em}%
}

% Shared chrome for CEI / student ethics PDFs — UDIT · ECSIT
% Built with XeLaTeX (CJK for 道). Input from documents in this directory.
\usepackage{graphicx}
\usepackage{fontspec}
\usepackage{hyperref}
\usepackage{xcolor}
\usepackage{tabularx}
\usepackage{fancyhdr}
\usepackage{calc}

% Single CJK glyph (道) via fontspec — no xeCJK dependency
\IfFontExistsTF{Songti SC}{%
  \newfontfamily\TTODdao{Songti SC}%
}{%
  \IfFontExistsTF{PingFang SC}{%
    \newfontfamily\TTODdao{PingFang SC}%
  }{%
    \IfFontExistsTF{STSong}{%
      \newfontfamily\TTODdao{STSong}%
    }{%
      \newcommand{\TTODdao}{}%
    }%
  }%
}

\definecolor{uditnavy}{HTML}{0B1F33}
\definecolor{uditink}{HTML}{1A1A1A}
\definecolor{rulegray}{HTML}{CCCCCC}

\newcommand{\assetspath}{../assets/}
\newcommand{\EcsitUrl}{https://www.udit.es/lineas-de-investigacion/grupo-de-investigacion-innovacion-y-tecnologia-desde-y-para-la-educacion-la-cultura-y-la-sociedad/}
\newcommand{\EcsitLink}{\href{\EcsitUrl}{udit.es $\rightarrow$ grupo ECSIT}}
\newcommand{\TTODfull}{{\TTODdao 道} The Tao of Development (TTOD)}

\newcommand{\LogoRow}{%
  \noindent
  \begin{tabularx}{\textwidth}{@{}lXr@{}}
    \includegraphics[height=0.95cm]{\assetspath udit-logo-print.png}
    & &
    \includegraphics[height=1.45cm]{\assetspath ecsit-logo-print.png}
  \end{tabularx}%
  \par\vspace{0.55em}
  \noindent{\color{rulegray}\hrule height 0.35pt}%
  \par\vspace{0.65em}
}

\newcommand{\LegalEntities}{%
  \noindent{\scriptsize\color{uditink}
  UDIT ESTUDIOS SUPERIORES INTERNACIONALES, S.L.\ · NIF B84014265 ·
  UDIT Escuela de Formación Profesional, S.L.\ · NIF B09963117.}%
}

% Role (left) | Name + email only (right) — no role text repeated
\newcommand{\AuthoritiesBlock}{%
  \noindent{\sffamily\bfseries\small Autoridades de investigación / Research authorities}\\[0.25em]
  \noindent
  \begin{tabularx}{\textwidth}{@{}>{\raggedright\arraybackslash}p{4.6cm}X@{}}
    \textbf{Coordinación grados tecnológicos} &
      Sandra Garrido ---
      \href{mailto:sandra.garrido@udit.es}{sandra.garrido@udit.es}\\[0.2em]
    \textbf{Dirección de grados} &
      Fernando Blázquez Piñeiro ---
      \href{mailto:fernando.blazquez@udit.es}{fernando.blazquez@udit.es}\\[0.2em]
    \textbf{IP del grupo ECSIT} &
      Rafael Conde Melguizo, PhD ---
      \href{mailto:rafael.conde@udit.es}{rafael.conde@udit.es}\\[0.2em]
    \textbf{CEI UDIT} &
      \href{mailto:comite.etica.investigacion@udit.es}{comite.etica.investigacion@udit.es}
  \end{tabularx}%
  \par\vspace{0.35em}%
  \LegalEntities
}

\newcommand{\StudyMetaPI}{%
  \noindent
  \begin{tabularx}{\textwidth}{@{}>{\raggedright\arraybackslash}p{4.6cm}X@{}}
    \textbf{Investigador principal} &
      Rubén Vega Balbás, PhD ---
      \href{mailto:ruben.vega@udit.es}{ruben.vega@udit.es}\\[0.15em]
    \textbf{Afiliación académica} &
      Grado Oficial en Desarrollo Full-Stack
      (\emph{La evolución lógica de la ingeniería informática hacia una
      titulación enfocada en el mundo real.})\\[0.15em]
    \textbf{Proyecto / producto} &
      \TTODfull\ · plataforma Oráculo · Cohorte Front-End~II\\[0.15em]
    \textbf{Grupo de investigación} &
      ECSIT · \EcsitLink\\[0.15em]
    \textbf{Comité} &
      \href{mailto:comite.etica.investigacion@udit.es}{comite.etica.investigacion@udit.es}\\[0.15em]
    \textbf{Código de estudio} &
      TTOD-FEII-2026-27\\[0.15em]
    \textbf{Versión del paquete} &
      2026-09-26 · cohorte pedagógica $n=8$
      (uso investigador: solo con consentimiento, $n\leq 8$)
  \end{tabularx}%
}

% Full Chicago (author–date) methods backbone — from docs/public/research/methodology.md
\newcommand{\ChicagoMethodsRefs}{%
\begin{itemize}
  \item Brown, Neil C.~C., and Mark Guzdial. 2024.
        ``Confidence vs Insight: Big and Rich Data in Computing Education Research.''
        In \emph{Proceedings of the 55th ACM Technical Symposium on Computer Science Education V.~1},
        158--64. \url{https://doi.org/10.1145/3626252.3630813}.
  \item Fischer, Björn, Berit Barthelmes, Sven Eric Panitz, Eva-Maria Iwer, and Ralf Dörner. 2024.
        ``Seeking Consent for Programming Process Data Collection with Trustee-Based Encryption.''
        In \emph{Proceedings of the 2024 ACM Conference on International Computing Education Research V.1},
        131--42. \url{https://doi.org/10.1145/3632620.3671125}.
  \item Heckman, Sarah, Jeffrey C.~Carver, Mark Sherriff, and Ahmed Al-Zubidy. 2022.
        ``A Systematic Literature Review of Empiricism and Norms of Reporting in Computing Education Research Literature.''
        \emph{ACM Transactions on Computing Education} 22 (1): Article~3.
        \url{https://doi.org/10.1145/3470652}.
  \item McGill, Monica M., Sarah Heckman, Christos Chytas, et~al. 2023.
        ``Conducting Sound, Equity-Enabling Computing Education Research.''
        In \emph{Proceedings of the 2023 Working Group Reports on Innovation and Technology in Computer Science Education},
        30--56. \url{https://doi.org/10.1145/3623762.3633495}.
  \item Runeson, Per, and Martin Höst. 2009.
        ``Guidelines for Conducting and Reporting Case Study Research in Software Engineering.''
        \emph{Empirical Software Engineering} 14 (2): 131--64.
        \url{https://doi.org/10.1007/s10664-008-9102-8}.
  \item Runeson, Per, Martin Höst, Austen Rainer, and Björn Regnell. 2012.
        \emph{Case Study Research in Software Engineering: Guidelines and Examples}.
        Hoboken, NJ: Wiley. \url{https://doi.org/10.1002/9781118181034}.
  \item Wohlin, Claes, Per Runeson, Martin Höst, Magnus C.~Ohlsson, Björn Regnell, and Anders Wesslén. 2012.
        \emph{Experimentation in Software Engineering}. Berlin: Springer.
        \url{https://doi.org/10.1007/978-3-642-29044-2}.
\end{itemize}%
}


\geometry{margin=2.2cm, top=2.0cm, bottom=2.2cm}
\setkomafont{section}{\sffamily\bfseries\large}
\setkomafont{subsection}{\sffamily\bfseries\normalsize}
\setkomafont{disposition}{\sffamily}
\setlength{\parskip}{0.32em}
\setlength{\parindent}{0pt}
\setstretch{1.04}

\hypersetup{
  colorlinks=true,
  urlcolor=uditnavy,
  linkcolor=uditnavy,
  pdftitle={Paquete institucional — autorización pedagógica y fondos TTOD},
  pdfauthor={Rubén Vega Balbás, PhD · UDIT / ECSIT},
  pdfsubject={Departamento / ECSIT / investigación — no es remisión CEI}
}

\pagestyle{fancy}
\fancyhf{}
\renewcommand{\headrulewidth}{0pt}
\fancyfoot[L]{\scriptsize\color{uditnavy} Institucional · \TTODfull}
\fancyfoot[C]{\scriptsize\color{uditink}\thepage}
\fancyfoot[R]{\scriptsize\color{uditink}v.\ 2026-09-26}

\newcommand{\AnnexTitle}[1]{%
  \clearpage
  \LogoRow
  \begin{center}{\Large\sffamily\bfseries #1}\end{center}
  \vspace{0.4em}
  \noindent{\color{rulegray}\hrule height 0.4pt}
  \vspace{0.5em}
}

\begin{document}

%----------------------------------------------------------------------
\LogoRow
\begin{center}
  {\LARGE\sffamily\bfseries Paquete institucional}\\[0.35em]
  {\large\sffamily Autorización pedagógica · Fondos de producción}\\[0.5em]
  {\normalsize\color{uditnavy} \TTODfull\ / Front-End~II}\\[0.35em]
  {\small Código TTOD-FEII-2026-27 · 26 septiembre 2026}
\end{center}

\vspace{0.5em}
\noindent{\sffamily\bfseries\small Destinatarios}\\[0.25em]
\noindent
\begin{tabularx}{\textwidth}{@{}>{\raggedright\arraybackslash}p{4.6cm}X@{}}
  \textbf{Coordinación grados tecnológicos} &
    Sandra Garrido ---
    \href{mailto:sandra.garrido@udit.es}{sandra.garrido@udit.es}\\[0.15em]
  \textbf{Dirección de grados} &
    Fernando Blázquez Piñeiro ---
    \href{mailto:fernando.blazquez@udit.es}{fernando.blazquez@udit.es}\\[0.15em]
  \textbf{IP del grupo ECSIT} &
    Rafael Conde Melguizo, PhD ---
    \href{mailto:rafael.conde@udit.es}{rafael.conde@udit.es}\\[0.15em]
  \textbf{Oficina investigación / OTRI} &
    Canal de fondos --- \emph{[completar email institucional]}
\end{tabularx}

\vspace{0.45em}
\noindent{\small\color{uditink}
\textbf{Este PDF no es la remisión al CEI.} El comité recibe el formulario Word +
\texttt{PAQUETE-REMISION-CEI-UDIT-ES.pdf}. Este pack pide (1)~confirmación pedagógica
firmable y (2)~apoyo presupuestario para producción $\geq 24$~meses.}

\vspace{0.4em}
\noindent{\sffamily\bfseries\small Contenido}
\begin{enumerate}
  \item Carta de autorización pedagógica (A1) --- Sandra y Fernando
  \item Solicitud de presupuesto de producción / hosting
  \item Separación de trámites (pedagogía / CEI / fondos)
\end{enumerate}

\vspace{0.35em}
\LegalEntities

%----------------------------------------------------------------------
\AnnexTitle{1.\ Solicitud de confirmación de encuadre pedagógico}

\noindent\textbf{De:} Rubén Vega Balbás, PhD ---
\href{mailto:ruben.vega@udit.es}{ruben.vega@udit.es}\\
\textbf{Asunto:} Confirmación de encuadre pedagógico --- grupo \TTODfull\ / Front-End~II

\subsection*{Solicitud}
Se solicita \textbf{confirmación escrita} (correo o esta hoja firmada) de que la
actividad del grupo \TTODfull\ / plataforma Oráculo, como entrega de
\textbf{Front-End~II}, queda \textbf{autorizada pedagógicamente} en el marco de los
grados tecnológicos de UDIT, con carga, calendario y evaluación coherentes con la
guía docente.

\subsection*{Qué se autoriza (docencia)}
\begin{enumerate}
  \item Pilotaje docente: construcción desde el \emph{cohort starter}, documentación
        de proceso, declaración de uso de IA y defensa.
  \item Coordinación de carga y coherencia de evaluación.
  \item Que la actividad funcione pedagógicamente \textbf{aunque no se recoja ningún
        dato} para investigación.
\end{enumerate}

\subsection*{Qué no se confunde}
El uso investigador seudonimizado se tramita ante el CEI
(\href{mailto:comite.etica.investigacion@udit.es}{comite.etica.investigacion@udit.es}).
Esta carta \textbf{no} autoriza investigación ni open data. Las calificaciones
\textbf{no} dependen del consentimiento investigador ni de la autorización de imagen.

\subsection*{Decisiones a marcar}
\renewcommand{\arraystretch}{1.3}
\begin{tabularx}{\textwidth}{@{}c>{\raggedright\arraybackslash}X c@{}}
  \toprule
  \textbf{\#} & \textbf{Decisión} & \textbf{Sí} \\
  \midrule
  1 & Autorizar el pilotaje docente TTOD Oráculo como entrega de Front-End~II & $\square$ \\
  2 & Confirmar coordinación de carga, calendario y coherencia de evaluación & $\square$ \\
  3 & Separar autorización pedagógica (esta carta) de la remisión al CEI & $\square$ \\
  \bottomrule
\end{tabularx}

\subsection*{Conformidad}
\vspace{0.6em}
\noindent\textbf{Coordinación grados tecnológicos --- Sandra Garrido}\\[0.35em]
Firma: \hrulefill\ Quad Fecha: \underline{\hspace{1.1cm}}~/~\underline{\hspace{1.1cm}}~/~\underline{\hspace{2cm}}\\
{\scriptsize\href{mailto:sandra.garrido@udit.es}{sandra.garrido@udit.es}}

\vspace{0.9em}
\noindent\textbf{Dirección de grados --- Fernando Blázquez Piñeiro}\\[0.35em]
Firma: \hrulefill\quad Fecha: \underline{\hspace{1.1cm}}~/~\underline{\hspace{1.1cm}}~/~\underline{\hspace{2cm}}\\
{\scriptsize\href{mailto:fernando.blazquez@udit.es}{fernando.blazquez@udit.es}}

\vspace{0.9em}
\noindent\textbf{Enterado (opcional) --- IP ECSIT · Rafael Conde Melguizo, PhD}\\[0.35em]
Firma: \hrulefill\quad Fecha: \underline{\hspace{1.1cm}}~/~\underline{\hspace{1.1cm}}~/~\underline{\hspace{2cm}}\\
{\scriptsize\href{mailto:rafael.conde@udit.es}{rafael.conde@udit.es}}

\vspace{0.5em}
\noindent{\scriptsize Contexto interno: \texttt{docs/research/DEPARTMENT-DECISION-BRIEF.md}.
Firmados fuera de git.}

%----------------------------------------------------------------------
\AnnexTitle{2.\ Solicitud de apoyo presupuestario --- producción $\geq 24$~meses}

\noindent Se solicita partida o apoyo para mantener \TTODfull\ (Oráculo)
\textbf{en producción} al menos \textbf{dos años}, de modo que el grupo disponga
de un producto vivo, el caso pedagógico (tras CEI) conserve referencia estable, y
ECSIT/UDIT cuenten con un demostrador de innovación docente.

\subsection*{Conceptos (importes a completar / validar con TI)}
\renewcommand{\arraystretch}{1.25}
\begin{tabularx}{\textwidth}{@{}c>{\raggedright\arraybackslash}X>{\raggedright\arraybackslash}p{3.6cm}@{}}
  \toprule
  \textbf{ID} & \textbf{Concepto} & \textbf{Est.~24 meses} \\
  \midrule
  H1 & Hosting / compute (VPS o imputación interna LAN/Lilith) & {[}completar €{]} \\
  H2 & Almacenamiento / copias de seguridad & {[}completar €{]} \\
  H3 & Dominio / TLS / DNS (si no absorbido por UDIT) & {[}completar €{]} \\
  H4 & Contingencia (10--15\,\%) & {[}completar €{]} \\
  \midrule
  & \textbf{Total propuesto} & \textbf{{[}completar €{]}} \\
  \bottomrule
\end{tabularx}

\vspace{0.35em}
\noindent{\scriptsize
Orden de magnitud orientativo de VPS pequeño: $\sim$15--60\,€/mes
($\sim$360--1\,440\,€ / 24 meses). Preferencia soberana: infra UDIT/LAN si es viable.}

\subsection*{Qué no se pide aquí}
Nómina del PI; LLM cloud; open data estudiantil; sustituir el dictamen del CEI.

\subsection*{Decisiones}
\renewcommand{\arraystretch}{1.3}
\begin{tabularx}{\textwidth}{@{}c>{\raggedright\arraybackslash}X c@{}}
  \toprule
  \textbf{\#} & \textbf{Decisión} & \textbf{Sí} \\
  \midrule
  1 & Apoyar continuidad en producción $\geq 24$ meses & $\square$ \\
  2 & Canal: grado / ECSIT / investigación / mixto (especificar al firmar) & $\square$ \\
  3 & Validar importes con servicios TI antes del compromiso & $\square$ \\
  \bottomrule
\end{tabularx}

\subsection*{Conformidad / enterado}
\vspace{0.55em}
\noindent\textbf{Sandra Garrido}\\
Firma: \hrulefill\quad Fecha: \underline{\hspace{1.1cm}}~/~\underline{\hspace{1.1cm}}~/~\underline{\hspace{2cm}}

\vspace{0.7em}
\noindent\textbf{Fernando Blázquez Piñeiro}\\
Firma: \hrulefill\quad Fecha: \underline{\hspace{1.1cm}}~/~\underline{\hspace{1.1cm}}~/~\underline{\hspace{2cm}}

\vspace{0.7em}
\noindent\textbf{Rafael Conde Melguizo, PhD (IP ECSIT)}\\
Firma: \hrulefill\quad Fecha: \underline{\hspace{1.1cm}}~/~\underline{\hspace{1.1cm}}~/~\underline{\hspace{2cm}}

\vspace{0.7em}
\noindent\textbf{Oficina de investigación / OTRI}\\
Nombre: \underline{\hspace{4.5cm}}\quad
Firma: \hrulefill\quad Fecha: \underline{\hspace{1.1cm}}~/~\underline{\hspace{1.1cm}}~/~\underline{\hspace{2cm}}

%----------------------------------------------------------------------
\AnnexTitle{3.\ Separación de trámites}

\renewcommand{\arraystretch}{1.25}
\begin{tabularx}{\textwidth}{@{}>{\raggedright\arraybackslash}p{3.2cm}X>{\raggedright\arraybackslash}p{3.8cm}@{}}
  \toprule
  \textbf{Trámite} & \textbf{Contenido} & \textbf{Documento} \\
  \midrule
  Pedagógico (A1) & Autorizar la actividad como entrega FE~II & Anexo~1 de este PDF \\
  Ética (CEI) & Uso investigador seudonimizado & Formulario Word + \texttt{PAQUETE-REMISION-CEI-UDIT-ES.pdf} \\
  Fondos & Producción $\geq 24$ meses & Anexo~2 de este PDF \\
  \bottomrule
\end{tabularx}

\vspace{0.6em}
\noindent Los tres trámites pueden avanzar en paralelo. Ninguno sustituye a los otros.

\vspace{1em}
\AuthoritiesBlock

\end{document}
